\section{Source audit, corrections of interpretation, and integration}
\label{audit:section}

\subsection{What the audit establishes}

The targeted review found no false theorem in the selected incoming
Tornheim statements or in the harmonic continuation used here. It
did find an attribution overlap that matters for presenting this
contribution: the canonical manuscript already gives the polynomial
Hurwitz transform and the complete-homogeneous coefficient mechanism
for all Stieltjes moments. The precise locations are
\src{chapters/07-integration.tex},
\src{stieltjes:sec:moments},
\src{stieltjes:thm:transform}, and
\src{stieltjes:thm:moment} \cite{source:manuscript}.
Accordingly, Corollary~\ref{harm:Stieltjesmoments} is classified as a
fully indexed arbitrary-endpoint restatement and extension, useful
for the new harmonic reduction, rather than a new general principle.

The formal difference-equation input, the Espinosa--Moll generalized
polygamma, and the classical Barnes identity are also credited
explicitly. Likewise, the boundary Gamma representation for $U$ in
\eqref{qtr:U} is an existing theorem specialized to the derivative
order needed here. These distinctions prevent an exact consequence
from being counted twice as a separate research discovery.

The reviewed source files were taken from the fixed commit stated in
Section~\ref{scope:section}. As an additional provenance check, the
44 inspected TeX files from the three newest incoming reports were
compared byte-for-byte with the corresponding ZIP contents in that
checkout; every comparison agreed. The packaged manifest records
the nine archive names, their Git blob identifiers, and SHA-256
hashes. An integration based on a later repository revision should
reconcile later changes explicitly.

\subsection{Current conjectures must retain their current labels}

The canonical conjectures
\src{cycloquot:conj:S6} and \src{s8new:conj:S8}
remain conjectural. The previously rejected $S_8$ coefficient vector
is a different candidate from the revised one. A rejection of the
former is not a rejection of the latter, and a numerical residual
for either candidate is not an analytic proof.

Two different computational statements in the incoming collection
also need to remain separate. The enlarged modular search reported
in \emph{Independent Orders} was interrupted before exhausting its
matrix; that record alone establishes neither membership nor
nonmembership. The later \emph{Exact Resonant Identities} report
establishes an exact obstruction for the complete Cayley ideal
together with its specifically enumerated additional row family.
That is a valid statement about the stated relation space. It does
not exclude additional distribution, double-shuffle, or functional
relations, and it is not a nonvanishing theorem for the actual
period \cite{qtr:Independent,ghp:ExactResonant}.

The analogous qualification is essential for the present
Theorem~\ref{qtr:kernel}: the formal dimension is proved under
symmetry, the $C$-axis identity, and the Euler-plane identity.
Its exactness within that system does not imply arithmetic
independence, nor does it claim to have imposed every possible
Tornheim functional equation.

The earlier source corrections concerning Tin's equation (22),
local polylogarithm estimates, and the rejected old $S_8$ vector
were already reported before this contribution. They are not
advertised here as newly discovered errors.

\subsection{Specific normalization rules supported by the new proofs}

The following are concrete instructions for using the formulas.
They also identify the corrections required if one extends the
existing calculus by a naive parameter substitution.

\paragraph{Retain the entire resonant amplitude.}
For the resolvent family the local term is
$r_p(z)A_p(s)x^{s-1}$, with
$A_p(s)=\prod_{j=1}^p(s-j)$. Its continued integral is
$r_pA_p(s)/s$. Subtracting only its residue at zero leaves
incorrect regular coefficients. The correct expression is
\eqref{res:completion}, with the entire polynomial amplitude
retained. At $p=1$, failure to do so misses the explicit constant
$1$ in \eqref{res:explicitcheck}. This is a proved correction
to that tempting extension, not an allegation that the pinned
source already contains the incorrect resolvent formula.

\paragraph{Keep the endpoint coordinate.}
All endpoint finite parts in the article use $x$, with the
cutoff specified in Section~\ref{res:section}. A nonlinear
endpoint substitution can change the finite constant of a
logarithmic divergence. The resolvent's spectral answer is
converted back to this stated coordinate before coefficient
extraction. Cauchy boundary values in \eqref{res:jump} are
asserted for the ordinary $p=0$ integral; no unproved transport
of divergent higher-derivative finite parts is implicit.

\paragraph{Do not differentiate a fixed index as though it were free.}
The deletion theorem is proved on the hyperplane where a
specified index equals $-m$. It permits every derivative in
the remaining independent free indices, because it is an
entire identity in those indices. It does not determine the
normal derivative transverse to the eliminated-index
hyperplane. That additional information is a separate research
problem.

\paragraph{Keep both Bernoulli endpoint conventions.}
The strict-gap deletion rule uses $B_1(0)=-1/2$ and
$B_1(1)=1/2$. Identifying them incorrectly deletes the last
term in
\[
 \mathcal H(s,0,t)=\mathcal H(s-1,t)
                 -\mathcal H(s,t-1)-\mathcal H(s,t).
\]
The verification suite deliberately corrupts that sign and
confirms that an independent finite nested sum detects it.

\paragraph{Restrict individual Tornheim rays before specialization.}
The individual fourth derivative in \eqref{qtr:quarticformula}
requires $(a+c)(b+c)\ne0$. A sum of cyclic germs can extend
across planes where individual terms are singular. The
holomorphic cyclic identity \eqref{qtr:J} has that extension;
its representation as a sum of individually restricted
derivatives is used only when every pairwise direction sum
is nonzero.

\paragraph{Specify what is centered and what is simultaneous.}
The involution in Theorem~\ref{ghp:thm:classification} changes
all even harmonic tails simultaneously. It does not require
each tail to occur to an even power separately. For example
$(\zeta(2)-H_n^{(2)})(\zeta(4)-H_n^{(4)})$ qualifies.
Conversely, a raw even-order harmonic number does not qualify:
$H_n^{(2)}$ produces the pole $-1/s$ at $s=0$.
This prevents two opposite misreadings of the classification.

\subsection{Proposed integration into the research collection}

The material separates naturally into four independently
readable additions: the centered generalized-harmonic
classification and sums; rational cotangent transforms and
balanced primitives; arbitrary-depth harmonic deletion and
the Lerch generator; and quartic and higher Tornheim jets.
Each has distinct label prefixes in the supplied TeX.

The incoming source questions can be annotated with the exact
answers in Table~\ref{scope:table}. Q11 receives a complete
classification. Q8 receives the stated rational sector.
Q9 receives the balanced primitive ladder for these kernels,
while a closed formula for every unbalanced Barnes product
is still open. The quartic directional question receives a
full normal form, and the all-order question receives an
exact answer for its specified formal relation system.

No existing conjecture should be promoted to a theorem
unless it is exactly one of these proved statements.
The package includes a claim ledger, proposed placement
notes, source hashes, and separate exact and numerical
verification records so that this distinction can survive
later merging and renumbering.
